% ===================================================================
% 3文書（研究計画書・説明文書・同意書）で共有する体裁設定
%
% 京都大学 医の倫理委員会「書類を作成する際のチェックリスト」の
%   ・使用されるフォントは「MSゴシック」「MS明朝」としてください
%   ・インデント、フォントサイズ、項目番号の表記の体裁を整えてください
% に対応するため、和文フォントと見出しの番号様式をこのファイルに集約する。
% kyodai-protocol.cls、explanation_template.tex、consent_template.tex の
% いずれからも読み込まれるため、3文書の体裁が自動的に揃う。
% ===================================================================

% -------------------------------------------------------------------
% 和文フォント: IPA明朝 / IPAゴシック
%
% チェックリストは「MS明朝」「MSゴシック」を指定するが、当該フォントは
% Microsoft専有でありビルド環境（texlive/texlive イメージ）に導入できない。
% 規定の趣旨である「Windowsとの互換性による文字化けの防止」はPDFへの
% フォント埋め込みによって満たしたうえで、書体としては指定に最も近い
% 無償の標準和文フォント（IPA明朝／IPAゴシック）を用いる。
% -------------------------------------------------------------------
\RequirePackage[ipa]{luatexja-preset}

% -------------------------------------------------------------------
% 見出しの番号様式
%
% 医の倫理委員会が示す「研究計画書に記載すべき事項」および実際に承認された
% 計画書で用いられている表記に揃える。
%   section       → 1.  2.  3. …
%   subsection    → 1)  2)  3) …
%   subsubsection → ①  ②  ③ …
%   paragraph     → 番号なし（太字見出し）
% -------------------------------------------------------------------
\RequirePackage{titlesec}

\newcommand{\jpcirclednum}[1]{%
  \ifcase#1\or ①\or ②\or ③\or ④\or ⑤\or ⑥\or ⑦\or ⑧\or ⑨\or ⑩%
  \or ⑪\or ⑫\or ⑬\or ⑭\or ⑮\or ⑯\or ⑰\or ⑱\or ⑲\or ⑳\else (\number#1)\fi}

\renewcommand{\thesection}{\arabic{section}.}
\renewcommand{\thesubsection}{\arabic{subsection})}
\renewcommand{\thesubsubsection}{\jpcirclednum{\value{subsubsection}}}

% 下位の見出しほど左を字下げし、階層を視覚的に示す
\titleformat{\section}[hang]{\normalfont\Large\bfseries}{\thesection}{0.7em}{}
\titlespacing*{\section}{0pt}{3.5ex plus 1ex minus .2ex}{2.3ex plus .2ex}
\titleformat{\subsection}[hang]{\normalfont\large\bfseries}{\thesubsection}{0.6em}{}
\titlespacing*{\subsection}{1em}{3.0ex plus 1ex minus .2ex}{1.8ex plus .2ex}
\titleformat{\subsubsection}[hang]{\normalfont\normalsize\bfseries}{\thesubsubsection}{0.4em}{}
\titlespacing*{\subsubsection}{2em}{2.6ex plus 1ex minus .2ex}{1.4ex plus .2ex}
\titleformat{\paragraph}[hang]{\normalfont\normalsize\bfseries}{}{0em}{}
\titlespacing*{\paragraph}{3em}{2.2ex plus 0.8ex minus .2ex}{1.0ex plus .2ex}

% 番号を付す深さはsubsubsection（①）まで。paragraphは番号なしの太字見出し。
\setcounter{secnumdepth}{3}
